Perkembangan \textit{Artificial Intelligence} (AI) dan \textit{Machine Learning} (ML) menuntut kompetensi baru pendidikan vokasional, khususnya Rekayasa Perangkat Lunak. Siswa SMK membutuhkan pengalaman praktis membangun solusi berbasis data sesuai standar industri digital. Mitra pengabdian ini, SMKN 2 Cimahi, menghadapi kendala berupa terbatasnya akses pelatihan praktis \textit{Machine Learning} dan \textit{Computer Vision}, belum optimalnya pembelajaran berbasis proyek AI, serta minimnya portofolio siswa untuk kesiapan kerja dan studi lanjut.

Sebagai solusi, diselenggarakan Workshop \textit{Machine Learning} untuk Prediksi Awal Kesehatan Tanaman melalui Analisis Citra Daun. Studi kasus ini dipilih karena bersifat visual, menggunakan dataset publik, tidak butuh perangkat keras khusus, dan relevan dengan kehidupan sehari-hari. Sistem dirancang sebagai alat \textit{screening} awal sehingga peserta memahami keterbatasan model ML dan pentingnya interpretasi prediksi secara kritis.

Kegiatan ini menggunakan pendekatan \textit{Project-Based Learning}, meliputi pengenalan konsep AI dan ML, eksplorasi dataset, pelatihan klasifikasi di \textit{Google Colab}, evaluasi model, hingga pembuatan aplikasi prediksi sederhana berbasis \textit{Streamlit}. Pelatihan didukung penuh oleh mahasiswa ULBI sebagai fasilitator teknis. Evaluasi keberhasilan program diukur melalui instrumen \textit{pre-test}, \textit{post-test}, penilaian \textit{mini project}, dan kuesioner kepuasan peserta.

Kegiatan ini diharapkan meningkatkan keterampilan praktis siswa dalam menerapkan \textit{Machine Learning} dan \textit{Computer Vision}, menghasilkan modul pembelajaran untuk sekolah, serta luaran publikasi di jurnal pengabdian nasional. Lebih lanjut, pengalaman belajar berbasis proyek ini akan memperkuat kompetensi digital dan melengkapi portofolio siswa SMK pada bidang teknologi berbasis data.

\vspace{0.5cm}
\noindent Kata Kunci : \textit{Machine Learning, Computer Vision, Project-Based Learning.}
